\section{Introduction}
\label{sec:introduction}

Time-correlation functions occupy a central position in statistical mechanics.
Through the Green--Kubo relations~\cite{green54a,kubo57a,kubo66a}, equilibrium fluctuations of a microscopic observable determine macroscopic transport coefficients; through the related Onsager reciprocal relations~\cite{onsager31a,onsager31b}, those coefficients exhibit the symmetries required by microscopic reversibility; and through van Hove correlation functions and their reciprocal-space transforms, they are tied to scattering experiments on dense fluids and soft matter~\cite{hansen90a}.
For essentially any observable derivable from the dynamics of a many-body system --- diffusion, ionic conductivity, shear viscosity, thermal conductivity, density fluctuations --- a correlation function can be written down, measured in simulation, and compared against experiment.

In practice, however, the information contained in these correlation functions is almost always compressed before it is used.
A viscosity is reported as a scalar; an ionic conductivity is reported as a scalar, sometimes decomposed into a Nernst--Einstein baseline and a distinct-correlation correction tailored to that specific coefficient~\cite{kashyap11a,fong20a,fong21a}; a dynamic structure factor is reported as a two-point function of wavevector and frequency.
Each of these reductions discards the particle-wise matrix structure from which it was assembled, and the self-versus-distinct decompositions that do exist are written down property-by-property rather than as a single framework that any transport coefficient of the relevant form can be dropped into.
Working backwards from such a coefficient to its microscopic definition recovers a sum over pairs of particles of an integrated velocity (or current) cross-correlation, and that object --- an $|\N| \times |\N|$ real symmetric matrix over all atoms in the system --- is seldom examined on its own terms.

This omission is not for lack of tools.
Correlation matrices constructed from high-dimensional data are studied extensively in other settings, where random matrix theory (RMT) provides a natural language for distinguishing physical structure from finite-sampling noise~\cite{TODO:plerou1999,TODO:laloux1999,TODO:bouchaud2009}.
In the molecular-dynamics literature, the closest relatives are \emph{configurational} in nature: Schlitter's formula~\cite{TODO:schlitter1993} and its refinements~\cite{TODO:andricioaei2001} compute an entropy from the spectrum of the mass-weighted positional covariance matrix, and essential-dynamics~\cite{TODO:amadei1993} / quasi-harmonic analyses diagonalise the same object to extract collective modes.
The matrix assembled from integrated velocity cross-correlations --- the object that actually appears in the Green--Kubo decomposition of transport coefficients --- has, to our knowledge, not been treated with the same spectral machinery.

In this work we do so, and in doing so make a contribution in two parts.
First, we give a species-generic Green--Kubo decomposition that expresses any transport coefficient of the form $\lambda \propto \int_{0}^{\infty} dt \, \langle A(t) \cdot A(0) \rangle$, with $A = \sum_{i} \eta_{\alpha(i)}\, a_{i}$ built from a single-particle primitive $a_{i}$ and a species-specific scalar $\eta_{\alpha}$, as a coarse linear functional of a single underlying object: the atom-wise time-integrated correlation matrix $C$.
This places self-diffusion, ionic conductivity, electrophoretic mobility, and related dissipative quantities on a common footing, rather than re-deriving a bespoke self / distinct decomposition for each one.
Second, we study $C$ directly.
We use its eigenvalue spectrum to identify physical correlation structure; we use its von Neumann entropy as a scalar summary of how correlation is distributed across the system; and we develop a Marchenko--Pastur-type null against which the bulk of the spectrum can be compared, so that physical signal is separable from finite-sample noise.
Using Lennard-Jones fluids, a controlled bond-addition protocol, water across its condensed phases, and a molten salt, we show that eigenvalue outliers of $C$ are diagnostic of bonding topology and that the von Neumann entropy of $C$ resolves first-order phase transitions at least as sharply as conventional thermodynamic observables, and in some cases more sharply.

The remainder of the paper is organised as follows.
Section~\ref{sec:theory} reviews the Onsager and Green--Kubo framework in a form suited to species-wise decomposition, introduces the matrix $C$ and the spectral quantities we use to characterise it, and states the RMT null.
Section~\ref{sec:methods} describes the simulation protocols, the construction of $C$ from a trajectory, and the noise-correction procedure.
Section~\ref{sec:experiments} presents the four experimental studies described above.
